% Generated by scripts/assemble.py; corrections belong in source/corrections.json.
\examhead{47}{上海大学 2026 年数学分析试卷(回忆版)}{49}
% p050-b04; source PDF page 50
\blockanchor{p050-b04}
(注. 题目来自考生回忆, 部分表述或数据可能有误)\par
% p050-b05; source PDF page 50
\blockanchor{p050-b05}
\question{1}求极限:\par
% p050-b06; source PDF page 50
\blockanchor{p050-b06}
(1) $\lim_{n\to \infty}\frac{7^n n!}{(3n)^n}.$\par
% p050-b07; source PDF page 50
\blockanchor{p050-b07}
(2)  $\lim_{n\to\infty}\left(\frac{1}{3}+\frac{2}{3^{2}}+\cdots+\frac{n}{3^{n}}\right).$\par
% p050-b08; source PDF page 50
\blockanchor{p050-b08}
\question{2}求积分 $\int_0^{\frac{\pi}{2}}\frac{1}{\sin x}\ln \frac{1 + a\sin x}{1 - a\sin x}\mathrm{d}x.$\par
% p050-b09; source PDF page 50
\blockanchor{p050-b09}
\question{3}求 $x^{2} + y^{2} = 2az$ 与 $(x^{2} + y^{2})^{2} = 2a^{2}xy$ 相交区域的体积.\par
% p050-b10; source PDF page 50
\par\Needspace{6\baselineskip}
\blockanchor{p050-b10}
\question{4}设\par
\nopagebreak[4]
% p050-b11; source PDF page 50
\blockanchor{p050-b11}
\displayformula{f (x, y) = (x - y) ^ {2} + \iint_ {\Sigma} x \left(z ^ {2} + e ^ {z}\right) \mathrm{d} y \mathrm{d} z + y \left(z ^ {2} + e ^ {z}\right) \mathrm{d} z \mathrm{d} x + \left(z f (x, y) - 2 e ^ {z}\right) \mathrm{d} x \mathrm{d} y.}
% p050-b12; source PDF page 50
\blockanchor{p050-b12}
求 $f(x,y)$ ，其中 $\Sigma :x^{2} + y^{2} + z^{2} = 1,z\geq 0,$ 方向取上侧.\par
% p050-b13; source PDF page 50
\blockanchor{p050-b13}
\question{5}设 $a_{n} \geq 0, a_{n+1} - a_{n} \leq \int_{n}^{n+1} \frac{1}{t \ln^{2} t + \sin^{2} t} \mathrm{d}t$, 判断 $\{a_{n}\}$ 是否收敛.\par
% p050-b14; source PDF page 50
\par\Needspace{6\baselineskip}
\blockanchor{p050-b14}
\question{6}设 $f(x)$ 在 $[a, b]$ 上非负且连续可导, 证明:\par
\nopagebreak[4]
% p050-b15; source PDF page 50
\blockanchor{p050-b15}
\displayformula{\left| \int_ {a} ^ {b} f ^ {3} (x) \mathrm{d} x - f ^ {2} (a) \int_ {a} ^ {b} f (x) \mathrm{d} x \right| \leq 2 \max _ {x \in [ a, b ]} | f ^ {\prime} (x) | \left(\int_ {a} ^ {b} f (x) \mathrm{d} x\right) ^ {2}}
% p050-b16; source PDF page 50
\blockanchor{p050-b16}
\question{7}设 $f(x)$ 在 $x = 0$ 的某邻域内连续可导，且 $\lim_{x \to 0} \frac{f(x)}{x} = 2$，判断级数 $\sum_{n=1}^{\infty} (-1)^{n+1} f\left(\frac{1}{\sqrt{n}}\right)$ 的敛散性.\par
% p050-b17; source PDF page 50
\par\Needspace{6\baselineskip}
\blockanchor{p050-b17}
\question{8}设 $D$ 是由 $y = x^{2}, y = 2x^{2}, xy = 1, xy = 4$ 围成的区域, 证明:\par
\nopagebreak[4]
% p050-b18; source PDF page 50
\blockanchor{p050-b18}
\displayformula{\int_ {\partial D} \frac {f (x y)}{x} \mathrm{d} x = - \frac {\ln 2}{3} \int_ {1} ^ {4} f ^ {\prime} (x) \mathrm{d} x.}
% p050-b19; source PDF page 50
\blockanchor{p050-b19}
\question{9}设函数列 $\{f_n(x)\}$ 在闭区间 $[a, b]$ 上点态收敛到 $f(x)$, 且 $f_n(x), f(x)$ 在 $[a, b]$ 上连续, $\{f_n(x)\}$ 关于 $n$ 单调, 证明: $\{f_n(x)\}$ 在 $[a, b]$ 上一致收敛于 $f(x)$.\par
% p050-b20; source PDF page 50
\blockanchor{p050-b20}
\question{10}设 $f(x, y) = \begin{cases} \frac{2xy^2}{x^2 + y^4}, & (x, y) \neq (0, 0); \\ 0, & (x, y) = (0, 0). \end{cases}$\par
% p050-b21; source PDF page 50
\blockanchor{p050-b21}
(1) 求 $f(x, y)$ 的偏导数.\par
% p050-b22; source PDF page 50
\blockanchor{p050-b22}
(2) 求 $f(x, y)$ 在 (0, 0) 处的方向导数.\par
% p050-b23; source PDF page 50
\blockanchor{p050-b23}
(3) $f(x,y)$ 在 $(0,0)$ 处是否可微？\par
% p050-b24; source PDF page 50
\blockanchor{p050-b24}
\question{11}解答如下问题:\par
% p050-b25; source PDF page 50
\blockanchor{p050-b25}
(1) 叙述反常积分 $\int_{a}^{+\infty} f(x, y) g(x, y) \mathrm{d}x$ 关于 $y$ 在区间 $I$ 上的 Abel 判别法.\par
% p050-b26; source PDF page 50
\blockanchor{p050-b26}
(2) $\int_0^{+\infty}\frac{\cos(2x)}{x + y} e^{-xy}\mathrm{d}x$ 在 $y\in [0,b]$ 上是否一致收敛？\par
